\documentclass[12pt,a4paper]{report}

% ------------------------------------------------------------
% PACKAGES
% ------------------------------------------------------------
\usepackage[utf8]{inputenc}
\usepackage[T1]{fontenc}
\usepackage{lmodern}
\usepackage{geometry}
\usepackage{setspace}
\usepackage{titlesec}
\usepackage{fancyhdr}
\usepackage{hyperref}
\usepackage{enumitem}
\usepackage{amsmath}
\usepackage{amsfonts}
\usepackage{relsize}

% ------------------------------------------------------------
% PAGE SETUP
% ------------------------------------------------------------
\geometry{margin=1in}
\setstretch{1.2}

% ------------------------------------------------------------
% HEADER / FOOTER
% ------------------------------------------------------------
\pagestyle{fancy}
\fancyhf{}
\fancyhead[L]{\leftmark}
\fancyhead[R]{\thepage}

% ------------------------------------------------------------
% QUESTION COUNTER
% ------------------------------------------------------------
\newcounter{question}[chapter]
\renewcommand{\thequestion}{\thechapter.\arabic{question}}

\newenvironment{question}[1][]{%
  \refstepcounter{question}%
  \vspace{1em}% space before each question
  \noindent\textbf{Question \thequestion.}%
  \ifx&#1&%
    % no optional argument
  \else
    \textbf{ #1}%
  \fi
  \quad
  \textit% italicize the question text
}{%
  \par\vspace{0.5em}% small space before the answer
}

% Optional "Answer" environment for clarity
\newenvironment{answer}{%
  \par\noindent\textbf{Answer:}\ }%
{\par\vspace{1em}}

% ------------------------------------------------------------
% TITLE SETTINGS
% ------------------------------------------------------------
\titleformat{\chapter}[display]
  {\normalfont\LARGE\bfseries}
  {\chaptertitlename\ \thechapter}{0.5em}{\huge}
\titlespacing*{\chapter}{0pt}{-5pt}{20pt}

% ------------------------------------------------------------
% DOCUMENT START
% ------------------------------------------------------------
\begin{document}

% Title Page
\begin{titlepage}
  \centering
  {\Huge \bfseries Answers to Lecture Notes \\[1em]}
  {\large MATH0029: Graph Theory and Combinatorics \\[0.5em]}
  {\large \today \\[2em]}
  \vfill
  {\large Author: Archie Worth\footnote{archie.worth.23@ucl.ac.uk}}
\end{titlepage}

% Table of Contents
\tableofcontents
\clearpage

% ------------------------------------------------------------
% INTRODUCTION (not numbered)
% ------------------------------------------------------------
\chapter*{Introduction}
\addcontentsline{toc}{chapter}{Introduction}

This document contains worked answers and notes corresponding to the lecture notes of MATH0029: Graph Theory and Combinatorics, as taught by Dr. António Girão in 2025, written by me (Archie Worth). Although they have been checked, expect about $\lceil n/2 \rceil$ to be actually correct.

\clearpage

\chapter{Basic Counting: Sets and Tuples}

\begin{question}
    How large is $[6]\choose{3}$?
\end{question}
\begin{answer}
  From Lemma 1.2(iv) we have that $|\binom{[6]}{3}| = \binom{6}{3} = 20$
\end{answer}

\begin{question} 
  For $A \subseteq \mathbb{N}$ define $S(A) = \Sigma_{a \in A}a$. For how many sets $A \in \binom{[6]}{3}$ is $S(A)$ odd? 
\end{question}
\begin{answer}
$\binom{[6]}{3} = \{ A \subseteq [6] : |A| = 3\}$. In order for $S(A)$ to be odd, we need an text odd number of odd elements (since odd + odd + even = even), and at least one odd number (as even + even + even = even). This then gives 2 possible configurations of sets: 
  \begin{align*}
    \{odd, odd, odd\} \text{ or } \{odd, even, even\} 
  \end{align*} 
  There are $\binom{3}{3}=1$ of the first type, and $\binom{3}{1} \binom{3}{2}= 3 \cdot 3 = 9$ of the second type, thus giving 10 total sets where $S(A)$ is odd. 
\end{answer}

\begin{question}
  Let $n \in \mathbb{N}$. How large are the following family of sets: 
  \begin{align*}
    \mathcal{F}_1 = \{ A \subseteq     &[n] : 1 \in A\} \\
    \mathcal{F}_3 = \{ A \subseteq [n] &: |A \cap [3]| \geq 2\} 
  \end{align*}
\end{question}

\begin{question}
  How many sets in $\binom{[6]}{3}$ do not contain two consecutive integers?
\end{question}

\begin{question}
  Let $n \in \mathbb{N}$ be odd. Find a simple expression for the sum
  \begin{align*}
    \sum\limits_{k>n/2}\binom{n}{k}
  \end{align*}
\end{question}

\begin{question}
  Let $n \in \mathbb{N}$ for which $k \in \mathbb{N}$ is $\binom{n}{k}$ maximal? Is $k$ unique? 
\end{question}


% ------------------------------------------------------------
% CHAPTER 2 EXAMPLE
% ------------------------------------------------------------
\chapter{Graphs}

\section{Basic Definitions}
\begin{question}
Prove that any graph of order $n \geq 2$ has two vertices of the same degree.    
\end{question}

\begin{question}
    Can you construct a 4-regular graph of order 7? If not, why not? How about a 3-regular graph of order 7?
\end{question}

\begin{question}
  Let $G$ be a graph of order $n \geq 2$ in which every vertex has degree at least $(n-1)/2$. Prove that for every pair of distinct non-adjacent vertices $u, v \in V(G)$ there is a third vertex $w$ such that $uw, vw \in E(G)$
\end{question}

\section{Examples of Graphs}
\begin{question}
    Let $n \in \mathbb{N}$. Which of the graphs $C_n, \, K_n,\,  K_{n,2n},\,  K_{n,n} \, \text{and } Q_n$ are regular? For those that are regular, calculate the degree of a vertex in the graph.
\end{question}

\begin{question}
  Let $n \in \mathbb{N}$. What are the sizes of the graphs: $C_n, \, K_n, \, K_{n,2n}, K_{n,n}, \, Q_n$?
\end{question}

\begin{question}
  Let $1 \leq a \leq b$ be integers. If $K_{a,b}$ has order $n$ what is the smallest and largest that the size of $K_{a,b}$ could be?
\end{question}

\section{Subgraphs and isomorphisms}

\begin{question}
  For which $a, b, t \in \mathbb{N}$ does $K_{a,b}$ contain a copy of $C_t$?  
\end{question}

\begin{question}
  How many copies of $K_{1,2}$ are there in $K_{a,b}$ (where $a, b \geq 2$)?
\end{question}

\begin{question}
  When does $K_{a,b}$ contain a cycle of length $a + b$?
\end{question}

\begin{question}
  How many copies of $C_4$ does $Q_n$ contain for $n \geq 2$?
\end{question}

\section{Walks, paths, and connectedness}


\section{Trees}
\section{Euler Circuits}
\section{Bipartite Graphs}
\section{Graph Colouring}

\chapter{The Probabilistic Method}
\section{Basics}
\section{Mycelskian Graphs}
\section{Random Graphs}
\section{Large Girth and Large Chromatic Number}

\chapter{Extremal Graph Theory}
\section{Hamilton Cycles}
\section{Forbidden Subgraphs: Mantel's Theorem}
\section{Forbidden Subgraphs: Turán's Theorem}
\end{document}
